\section{Results}
\label{sec:results}

\subsection{In the matched comparison, the discrimination AUC fell within the equivalence band}
\label{sec:res-main}

The holdout gave 4,554 stalls in the no-event stratum, each with up to five controls, for 709,352 pairs among 139
writers (effective number of writers 50.5). For both signals and both tolerances, the 95\% interval lay inside the band
(Table~\ref{tab:main}, Figure~\ref{fig:forest}), so the preregistered test was passed [C]. The share of tied pairs was
6.5\% for the Stuck-state probability and 0.34\% for the friction axis at $\varepsilon = 0.001$, far from the
degeneracy rule. The writer-equal version gave 0.491 and 0.488 [C-D]. An audit of how the two estimate rows of each change were picked is in
Appendix~\ref{app:definitions} [E-H].

\begin{table}
  \caption{Discrimination AUC of the 30-second change at $u$ in the matched comparison, CS1 holdout, no-event stratum
  (4,554 stalls, 709,352 pairs, 139 writers). [C] except the last column [C-D].}
  \label{tab:main}
  \centering\small
  \begin{tabularx}{\linewidth}{>{\raggedright\arraybackslash}X c c c l c}
    \toprule
    Signal & $\varepsilon$ & AUC & 95\% interval & Result & Writer-equal \\
    \midrule
    Change in the Stuck-state probability, $\Delta p_{\mathrm{stuck}}$ & 0.001 & 0.489 & 0.479--0.4995 & equivalent & 0.491 \\
     & 0 & 0.489 & 0.479--0.5001 & equivalent & 0.492 \\
    Change in the smoothed friction axis, $\Delta x$ & 0.001 & 0.487 & 0.476--0.497 & equivalent & 0.488 \\
     & 0 & 0.487 & 0.476--0.497 & equivalent & 0.488 \\
    \bottomrule
  \end{tabularx}

  \smallskip
  \parbox{\linewidth}{\footnotesize Note. All four point estimates lie slightly below 0.5; three of the four intervals
  exclude 0.5, so these three are both statistically different from 0.5 and equivalent \cite{lakens2018equivalence}.}
\end{table}

\begin{figure}
  \centering
  \Description{A dot plot with horizontal intervals. Each row is an analysis; for each row two points with 95\% intervals
  are shown, one for the change in the Stuck-state probability and one for the change in the friction axis. A grey
  vertical band marks the equivalence band from 0.45 to 0.55. The main analysis lies at about 0.49 for both signals,
  inside the band. The writer-equal version, the not-at-floor pairs selected at the start point, comparisons within
  assignments, matching on whitespace or symbol, the version that drops stalls with events inside the pause, the event
  stratum, three-way matching on position, the 2021 CS1 writers and the essay data all lie inside the band, between
  about 0.47 and 0.52. Only the not-at-floor pairs selected at the end point lie lower, at about 0.465, with intervals
  crossing 0.45.}
  \begin{tikzpicture}
    \begin{axis}[
      width=0.66\linewidth, height=11cm,
      xmin=0.42, xmax=0.58, xtick={0.42,0.45,0.50,0.55,0.58},
      x tick label style={/pgf/number format/fixed, /pgf/number format/precision=2, font=\footnotesize},
      xlabel={Discrimination AUC ($\varepsilon = 0.001$) with 95\% interval}, label style={font=\footnotesize},
      ytick={1,...,11}, y dir=reverse, ymin=0.4, ymax=11.6,
      yticklabels={{Main [C]}, {Writer-equal [C-D]}, {Not at floor, start [C-D]}, {Within assignment [C-D]},
        {Whitespace or symbol [C-D]}, {No events in pause [C-D]}, {Event stratum [C-D]}, {Not at floor, end [C-D]},
        {Three-way position [E-H]}, {CS1 2021, 31 writers [E]}, {KLiCKe essays [D]}},
      y tick label style={font=\footnotesize, align=right},
      legend style={font=\footnotesize, draw=none, at={(0.5,1.01)}, anchor=south, legend columns=1}, legend cell align=left,
      xmajorgrids, major grid style={black!12},
    ]
      \fill[figband] (axis cs:0.45,0.4) rectangle (axis cs:0.55,11.6);
      \draw[black!60, dashed] (axis cs:0.5,0.4) -- (axis cs:0.5,11.6);
      \addplot+[only marks, mark=*, color=figc1, error bars/.cd, x dir=both, x explicit]
        coordinates {
          (0.4892,0.85) -= (0.0105,0) += (0.0103,0)
          (0.4912,1.85) -= (0.0182,0) += (0.0178,0)
          (0.4979,2.85) -= (0.0215,0) += (0.0206,0)
          (0.4951,3.85) -= (0.0113,0) += (0.0107,0)
          (0.4901,4.85) -= (0.0104,0) += (0.0111,0)
          (0.5047,5.85) -= (0.0121,0) += (0.0135,0)
          (0.4946,6.85) -= (0.0177,0) += (0.0204,0)
          (0.4649,7.85) -= (0.0218,0) += (0.0207,0)
          (0.4898,8.85) -= (0.0108,0) += (0.0108,0)
          (0.4925,9.85) -= (0.0079,0) += (0.0083,0)
          (0.506,10.85) -= (0.005,0) += (0.005,0)
        };
      \addlegendentry{Change in the Stuck-state probability}
      \addplot+[only marks, mark=square*, color=figc2, error bars/.cd, x dir=both, x explicit]
        coordinates {
          (0.4867,1.15) -= (0.0109,0) += (0.0108,0)
          (0.4875,2.15) -= (0.0156,0) += (0.0173,0)
          (0.4977,3.15) -= (0.0200,0) += (0.0190,0)
          (0.4920,4.15) -= (0.0112,0) += (0.0115,0)
          (0.4890,5.15) -= (0.0115,0) += (0.0113,0)
          (0.4999,6.15) -= (0.0163,0) += (0.0151,0)
          (0.4960,7.15) -= (0.0189,0) += (0.0202,0)
          (0.4656,8.15) -= (0.0173,0) += (0.0162,0)
          (0.4877,9.15) -= (0.0120,0) += (0.0122,0)
          (0.4861,10.15) -= (0.0107,0) += (0.0129,0)
          (0.513,11.15) -= (0.005,0) += (0.005,0)
        };
      \addlegendentry{Change in the smoothed friction axis}
    \end{axis}
  \end{tikzpicture}
  \caption{Discrimination AUC of the main analysis, the sensitivity analyses and the exploration, with the equivalence
  band shaded. Labels give the status of each analysis (Table~\ref{tab:status}). Values with $\varepsilon = 0$ are
  in Table~\ref{tab:main} and Appendix~\ref{app:sensitivity}.}
  \label{fig:forest}
\end{figure}

The levels of the signals at $u$ discriminated slightly more than their changes, and also stayed inside the band:
0.524 (0.510--0.537) for the Stuck-state probability and 0.521 (0.508--0.535) for the friction axis [C-D]. In the essay
data seen before the study, the same statistic gave 0.506 (0.501--0.511) and 0.513 (0.508--0.518) for the changes and
0.569 and 0.580 for the levels [D]. With the 90th percentile of each writer's controls as the threshold, 10.3\% of stall
run-ups exceeded it on the Stuck-state probability and 9.3\% on the friction axis, against 9.6\% of controls [C-D] (the control rate is about 10\% by construction); the differences, +0.7 and
$-0.3$ points, had writer-resampled intervals of $-0.2$ to 1.6 and $-1.2$ to 0.5 points [E-H], both including 0. These
are rates in the selected set of stalls and controls, not the performance of a trigger with a threshold set
independently and run continuously.

\subsection{On the same pairs, learners on simple quantities of preceding typing exceeded the two changes outside the registration}
\label{sec:res-learners}

This comparison sets the estimator's signals beside standard quantities judged against the same partner. In the
registered description [C-D], learners on 18 quantities of the preceding typing (Section~\ref{sec:design}) passed their
gate on the holdout: with a known difference added to one quantity so that its own AUC was 0.580, both recovered an
interval with a lower bound above 0.55. Logistic regression reached 0.567 (0.5497--0.578) and gradient boosting 0.540
(0.525--0.558), with intervals from 200 resamples of writers; both intervals cross 0.55, so the learners neither
established nor excluded an AUC above 0.55, for logistic regression only at the resolution of 200 resamples. Applied to
unmatched windows, the same learners gave 0.529 (0.512--0.546) and 0.526 (0.511--0.541), so they did not exceed a lower
bound of 0.55 there either (Appendix~\ref{app:learners}).

Outside the registration [E-H], with the learners' statistic ($\varepsilon = 0$) and over their 200 resamples (which
reproduced the registered values exactly), logistic regression exceeded the two 30-second changes by 0.078 (95\% interval
0.057--0.093) and 0.081 (0.057--0.097), and the levels of the signals at $u$ by 0.042 (0.019--0.059) and 0.046
(0.021--0.064); gradient boosting exceeded the changes by 0.051 and 0.054 but not clearly the levels. The learners can
use either direction of a quantity, whereas each signal was ranked in a fixed direction; read without direction (the
larger of the AUC and one minus it for the signals), logistic regression exceeded the changes by 0.057 (0.032--0.070)
and 0.054 (0.033--0.066), and gradient boosting by 0.030 and 0.027 (lower ends 0.009 and 0.008). Taken alone, the
quantities closest to the learners were levels over the 30~s before $u$: stalls were preceded by fewer key presses (AUC
0.446) and fewer characters per minute (0.443), whereas every 30-second change among the 18 quantities lay between
0.492 and 0.508; we did not decompose the learners. Because condition A becomes true earlier after less net production
(Appendix~\ref{app:definitions}), the stall definition could itself tie stalls to less preceding typing. Among the 3,426
stalls whose onset it did not bring forward ($s - u = 30$~s), the AUC of the number of key presses was 0.467 (0.393
among the 1,057 with $s - u < 30$~s; the other 71 stalls, with $s - u > 30$~s, are in neither group), and the
comparison, rerun on these stalls with the learners refitted, gave logistic regression 0.550 (0.527--0.562); read
without direction, it exceeded the two changes by 0.037 (0.012--0.056) and 0.031 (0.009--0.044), while gradient
boosting did not clearly exceed them (0.016 and 0.010, intervals including 0). By the reading committed before this
computation, the difference is therefore not explained by the definition of a stall alone, although it is smaller
there (Appendix~\ref{app:learners}).

\subsection{The main sensitivity analyses were also equivalent}
\label{sec:res-sens}

Of the preregistered sensitivity analyses, four bear on how the main result should be read [C-D]
(Figure~\ref{fig:forest}). The Stuck-state probability is mostly at the floor of the estimate: 61\% of stalls and 65\%
of controls had a probability below 0.02 at $u$, and none was above 0.98. Keeping only the pairs whose stall and control were both off the floor
at the start of the window, $u - 30$~s, left 2,090 stalls and gave 0.498 for both signals, equivalent under both
tolerances. Drawing controls only within the same assignment gave 0.495 and 0.492; matching on whether the character at
$u$ is whitespace or a symbol, instead of a newline, gave 0.490 and 0.489. Extending the stratum window to the end of
the pause, which removes stalls with a run, paste or switch inside the pause, left 2,427 stalls and gave 0.505
(0.493--0.518) and 0.500 (0.484--0.515). All four were equivalent for both signals and both tolerances. Selecting
off-floor pairs at the end point $u$ instead gave about 0.465 with intervals crossing 0.45 (undecidable on the lower
side); that version selects on the level reached at $u$ and thus on the change itself, which is why the start-point
version was registered as the main one. The event stratum gave 0.495 and 0.496 (effective writers 35.7). The remaining
sensitivity analyses, all equivalent, and the spread over the seed of the control draw are in
Appendix~\ref{app:sensitivity}.

\subsection{The unmatched comparison did not show a rise beyond the band}

In the secondary comparison of windows within writers, the unmatched AUC was 0.497 (0.491--0.503) for the Stuck-state
probability and 0.494 (0.489--0.4997) for the friction axis [C-D]. The registered condition, a lower bound above 0.55,
was not met, so we report only the equivalence in the matched comparison and do not contrast the two
(Appendix~\ref{app:unmatched}).

\subsection{Outside the registration, the estimator barely registered changes of the observed size}
\label{sec:res-bounds}

Three explorations outside the registration, each with conditions committed before it was computed, bound how the
equivalence can be read. They do not enter the decision.

\paragraph{Positive control of the estimator [E].} On the 319 seen writers, we rebuilt the estimate rows after changing
the key presses in the 30~s before every stall's $u$, keeping the stalls, controls and draws fixed. The intervals
between key presses were multiplied by $k$ (keeping the press at $u$), or a backspace and a retyped key were inserted
into a share $r$ of the intervals (Table~\ref{tab:poscontrol}). Because the estimator carries its state forward, the change also altered 34--47\% of the drawn
controls at the larger sizes, so we also report a version that keeps the controls at their unchanged values. The
observed difference was small and ran the other way: in the 30~s before $u$, the median of the mean interval between
presses was 598~ms for stalls and 605~ms for controls (a ratio of about 1.01), and the median share of backspaces 0.091
and 0.100. The injected changes ran opposite to this difference, and changes in its direction were not made. Slowing the run-up
by 2--10\%, more than the difference, hardly moved the AUC: 0.492--0.499. With controls
fixed, the friction axis left the band only when the intervals were doubled or half of them received a correction; the
Stuck-state probability did not leave the band even when the intervals were tripled (0.553, an interval crossing 0.55)
and fell when corrections were added. The levels of the signals at $u$, 0.526 and 0.527 unchanged, left the band at
smaller changes than their 30-second changes (Table~\ref{tab:poscontrol-all}).

\begin{table}
  \caption{Positive control on the seen CS1 writers [E]: discrimination AUC ($\varepsilon = 0.001$) after changing the key
  presses in the 30~s before each stall's $u$, with the controls re-estimated or kept at their unchanged values. Values
  whose 95\% interval lies outside the band are in bold; the half-width of the intervals was 0.007--0.008. Selected
  sizes; all sizes and the levels are in Table~\ref{tab:poscontrol-all}.}
  \label{tab:poscontrol}
  \centering\small
  \begin{tabular}{l cc cc}
    \toprule
    & \multicolumn{2}{c}{Controls re-estimated} & \multicolumn{2}{c}{Controls fixed} \\
    Change & $\Delta p_{\mathrm{stuck}}$ & $\Delta x$ & $\Delta p_{\mathrm{stuck}}$ & $\Delta x$ \\
    \midrule
    None & 0.492 & 0.490 & 0.492 & 0.490 \\
    Intervals $\times 1.02$ & 0.493 & 0.493 & 0.493 & 0.492 \\
    Intervals $\times 1.1$ & 0.499 & 0.498 & 0.498 & 0.496 \\
    Intervals $\times 1.5$ & 0.518 & 0.524 & 0.515 & 0.523 \\
    Intervals $\times 2$ & 0.537 & \textbf{0.564} & 0.532 & \textbf{0.564} \\
    Intervals $\times 3$ & \textbf{0.562} & \textbf{0.627} & 0.553 & \textbf{0.626} \\
    Corrections, $r = 0.05$ & 0.496 & 0.512 & 0.490 & 0.509 \\
    Corrections, $r = 0.25$ & 0.500 & \textbf{0.562} & 0.480 & 0.557 \\
    Corrections, $r = 0.5$ & 0.495 & \textbf{0.563} & 0.469 & \textbf{0.558} \\
    \bottomrule
  \end{tabular}
\end{table}

\paragraph{Scaled inputs at their floor or ceiling [E-H].} On the holdout's main stratum we replayed every assignment
again to read the estimator's window quantities. At
$u$, each of the four scaled inputs was pinned at 0 or 1 in a large share of stalls: the flight time at its floor in
72\% (at its ceiling in 10\%), the share of pauses after deletions in 83\% (9\%), the correction ratio in 32\% (35\%) and
the number of pauses in 54\% (17\%); the number of pauses moved from its ceiling at $u - 30$~s (51\%) to its floor at
$u$. The 30-second change was exactly zero in 63\% of stalls for the flight time and 73\% for the share of pauses after
deletions. Controls were alike (Table~\ref{tab:phi}). Part of the equivalence may therefore come from the floor and
ceiling of the scaling on this input.

\paragraph{Later-pause controls [E-H].} A trigger that fires after a pause has lasted $x$ seconds compares stalls with
pauses that outlasted $x$. Restricting the controls in the holdout to pauses of at least 5~s or 10~s gave 0.491 and 0.491,
and 0.494 and 0.497, all inside the band (effective writers 49.8 and 48.6).

\subsection{In the exploration outside the registration, the discrimination AUC also stayed inside the band}
\label{sec:res-explore}

After the confirmation we applied the same procedure, with conditions committed before each computation, to other data
(Table~\ref{tab:explore}). In the 31 seen writers of the 2021 CS1 edition, all four combinations were inside the band
[E]. Matching stalls and controls on three positions at once (after a newline, after a token boundary, within a token)
in the holdout gave 0.490 and 0.488 [E-H]; stalls were slightly more frequent after a token boundary than controls (53.9\%
against 50.9\%). The procedure also ran on the full essay corpus and on the seen co-writing data (Appendix~\ref{app:explore}).

\begin{table}
  \caption{Exploration outside the registration: discrimination AUC of the 30-second change ($\varepsilon = 0.001$; the
  value with $\varepsilon = 0$ differs by at most 0.001).}
  \label{tab:explore}
  \centering\small
  \begin{tabularx}{\linewidth}{>{\raggedright\arraybackslash}X l c c}
    \toprule
    Data & Status & $\Delta p_{\mathrm{stuck}}$ & $\Delta x$ \\
    \midrule
    CS1 2021 edition, 31 seen writers & [E] & 0.493 (0.485--0.501) & 0.486 (0.475--0.499) \\
    CS1 2019 holdout, three-way position & [E-H] & 0.490 (0.479--0.5006) & 0.488 (0.476--0.4999) \\
    \bottomrule
  \end{tabularx}
\end{table}
